\chapter{2-Dimensional Cellular Automata}

\section{Preliminaries}

A 2-dimensional lattice is an array of cells tiling an infinite planar
surface, each taking a value from a finite set at each tick of the
clock; the state at time $t$ is the doubly infinite array
$C^{(t)} = \bigl[\,C^{(t)}_{ij}\,\bigr]$.  As before we consider only
periodic boundary conditions, but now the wrap-around occurs in two
directions: right wraps to left and bottom wraps to top.  Fixing $n$
rows and $m$ columns, the state becomes the finite array
\begin{equation}
  C^{(t)} =
  \begin{bmatrix}
    C^{(t)}_{0,0} & C^{(t)}_{0,1} & \cdots & C^{(t)}_{0,m-1}\\
    C^{(t)}_{1,0} & C^{(t)}_{1,1} & \cdots & C^{(t)}_{1,m-1}\\
    \vdots & \vdots & \ddots & \vdots\\
    C^{(t)}_{n-1,0} & C^{(t)}_{n-1,1} & \cdots & C^{(t)}_{n-1,m-1}
  \end{bmatrix}.
\end{equation}
One way to visualize the periodic boundary condition: join the top and
bottom of the array to form a tube, then join the ends of the tube.
Our lattice has the topology of a \emph{torus}.  Equivalently, tile the
infinite plane with copies of the $n\times m$ array, requiring
\begin{equation}
  a \equiv b \pmod n,\ \ c \equiv d \pmod m
  \;\Longrightarrow\; C^{(t)}_{a,c} = C^{(t)}_{b,d}:
\end{equation}
all arithmetic on the first index is done modulo $n$, on the second
modulo $m$.

As before, we consider homogeneous CA (the same relative neighborhood
for every cell), let $k$ denote the number of states, identified with
$0,1,\dots,k-1$, and compute cell values modulo $k$.  The number of
possible states of such a CA is $k^{nm}$; in
Section~\ref{sec:vonneumann}, with $k=256$
and $n=m=500$, the number of states is $256^{500\cdot500} >
10^{602059}$.

Before the examples we discuss the generalization of the convolution
kernels found so useful in the previous chapter.  These 2-dimensional
kernels are used extensively in image processing.

\section{2-Dimensional Convolution Kernels}

One of the most frequently used procedures in image processing is
filtering an image by convolving it with a 2-dimensional kernel.  The
procedure closely parallels the 1-dimensional case.  For example, the
kernel
\[
  \begin{bmatrix*}[r] 0 & -1 & 0\\ -1 & 4 & -1\\ 0 & -1 & 0
  \end{bmatrix*}
\]
applied to a digital image serves as an \emph{edge extractor}, since it
maps homogeneous regions to zero.  Center the kernel over each pixel
and form the multiply-and-add.  For the image fragment
\[
  \begin{bmatrix} 1 & 2 & 3\\ 5 & \mathbf{6} & 7\\ 8 & 9 & 1
  \end{bmatrix}
\]
the output value at the central pixel (input value 6) is
\[
  0\cdot1 + (-1)\cdot2 + 0\cdot3
  + (-1)\cdot5 + 4\cdot6 + (-1)\cdot7
  + 0\cdot8 + (-1)\cdot9 + 0\cdot1 = 1 .
\]
In general, for an $N\times M$ image $x_{ij}$ and a
$(2Q\!-\!1)\times(2Q\!-\!1)$ kernel $k_{fh}$, the output with periodic
boundaries is
\begin{equation}
  y_{ij} \;=\; \sum_{f=1-Q}^{Q-1}\ \sum_{h=1-Q}^{Q-1}
     k_{fh}\; x_{(i+f)\bmod N,\ (j+h)\bmod M}.
\end{equation}
Additive 2-dimensional CA are amenable to detailed mathematical
analysis and admit a polynomial representation similar to the
1-dimensional case, but with two variables
(Section~\ref{sec:poly2d}).

\section{Interlude: The Same Kernels as Image Filters}
\label{sec:imageproc}

Before returning to CA proper it is worth pausing to actually
\emph{see} the image-processing use of this machinery, since the two
computations are literally identical: one generation of an additive CA
is one convolution pass over an image, minus the mod-$k$ reduction.
Figure~\ref{fig:imageproc} applies two kernels from this chapter to a
synthetic grayscale test image (a gradient background bearing a disk, a
square, and a triangle, with mild noise).

The center panel shows a single pass of the edge-extracting Laplacian
kernel introduced in the previous section: homogeneous regions ---
where a pixel equals the average
of its four neighbors --- are mapped to zero, so only the boundaries of
the shapes survive.  The right panel shows six iterated passes of the
smoothing kernel
\[
  \tfrac{1}{16}
  \begin{bmatrix} 1&2&1\\ 2&4&2\\ 1&2&1 \end{bmatrix},
\]
the outer product $[1\;2\;1]^{\mathsf T}[1\;2\;1]$ of the very kernel
used in the worked example of Section~3.1, normalized to sum to 1.
Iterating a smoothing convolution is running a CA whose long-run
behavior is the diffusion (heat) equation --- a first concrete instance
of the claim of Chapter~1 that CA yield discrete approximations to
differential equations.  Where an image processor stops after one or a
few passes, a CA keeps going; the difference between the two fields is,
as promised, one of intent rather than mechanism.

\begin{figure}[htbp]\centering
\includegraphics[width=\textwidth]{ch4_image_processing}
\caption{The kernels of this chapter doing their day job.  Left: a
synthetic test image.  Center: one pass of the Laplacian
edge-extraction kernel (absolute value shown).  Right: six passes of
the normalized smoothing kernel --- iterated convolution as discrete
diffusion.}
\label{fig:imageproc}
\end{figure}

\section{Example: $k=256$, $n=m=500$, the von Neumann kernel}
\label{sec:vonneumann}

For this example the state at generation $t$ is a $500\times500$ array
with entries in $\{0,\dots,255\}$; wrap-around is imposed by modulo-500
arithmetic on the indices.  The neighborhood of a cell is the cell
itself with its left, right, upper, and lower neighbors, and the rule
is the neighborhood sum modulo 256:
\begin{equation}
  C^{(t+1)}_{ij} = C^{(t)}_{i,j-1} + C^{(t)}_{i-1,j} + C^{(t)}_{ij}
  + C^{(t)}_{i,j+1} + C^{(t)}_{i+1,j} \pmod{256},
\end{equation}
compactly expressed by the kernel
\begin{equation}
  \begin{bmatrix} 0&1&0\\ 1&1&1\\ 0&1&0 \end{bmatrix}.
\end{equation}
We start from the simplest initial state: all cells zero except one
cell of value 1.  In the following images black represents zero and
colors the other 255 states.

From this simple seed an ever-growing mandala of intricate complexity
and symmetry unfolds (growing up to the limiting $500\times500$
dimensions of the space).  Figure~\ref{fig:vn} shows the seed after 100
and after 250 generations.  Figure~\ref{fig:vnseeds} shows the descent
of a different initial state: 256 seeds with values $0,\dots,255$
arranged on a grid, after 13 generations --- each seed the germ of its
own distinct mandala.  (The seed of value 0 is a perfectly valid seed;
it simply proves infertile, leaving a conspicuous gap in the garden.)

\begin{figure}[htbp]\centering
\includegraphics[width=0.48\textwidth]{ch4_vonneumann_gen100}\;
\includegraphics[width=0.48\textwidth]{ch4_vonneumann_gen250}
\caption{The von Neumann--kernel CA grown from a single seed:
generation 100 (detail, left) and generation 250 (right).}
\label{fig:vn}
\end{figure}

\begin{figure}[htbp]\centering
\includegraphics[width=0.7\textwidth]{ch4_vonneumann_seeds_gen13}
\caption{Generation 13 from an initial state of 256 seeds with values
$0,\dots,255$; the value-0 seed, though valid, is infertile.}
\label{fig:vnseeds}
\end{figure}

\section{Example: $k=256$, $n=m=500$, a signed kernel}
\label{sec:kernel43}

With the same lattice and state count as the previous example, we now
take the neighborhood to be the cell together with its eight nearest
neighbors, and the transition rule given by the kernel
\begin{equation}
  \begin{bmatrix*}[r] -1 & 2 & -1\\ 2 & -4 & 2\\ -1 & 2 & -1
  \end{bmatrix*},
\end{equation}
that is,
\begin{equation}
\begin{aligned}
  C^{(t+1)}_{ij} ={}& -C^{(t)}_{i-1,j-1} + 2C^{(t)}_{i-1,j}
                     - C^{(t)}_{i-1,j+1}
                     + 2C^{(t)}_{i,j-1} - 4C^{(t)}_{ij}
                     + 2C^{(t)}_{i,j+1}\\
                   & - C^{(t)}_{i+1,j-1} + 2C^{(t)}_{i+1,j}
                     - C^{(t)}_{i+1,j+1} \pmod{256}.
\end{aligned}
\end{equation}
Starting again from a single seed of value 1 (white now representing
zero), we once more see a mandala of intricate complexity and symmetry
unfold; Figure~\ref{fig:k43} shows generation 250 and a magnified
detail.

\begin{figure}[htbp]\centering
\includegraphics[width=0.48\textwidth]{ch4_kernel43_gen250}\;
\includegraphics[width=0.48\textwidth]{ch4_kernel43_gen250_detail}
\caption{The signed-kernel CA of Section~\ref{sec:kernel43} at
generation 250 (left) with a magnified central detail (right).}
\label{fig:k43}
\end{figure}

\section{Polynomial Representation of 2-Dimensional Additive CA}
\label{sec:poly2d}

Two-dimensional additive CA have polynomial representations paralleling
Section~\ref{sec:poly1d}.  A state --- an $n\times m$ matrix
$[a_{ij}]$ of integers mod $k$ --- is represented by the two-variable
characteristic polynomial
\begin{equation}
  A(x,y) \;=\; \sum_{i=0}^{n-1}\sum_{j=0}^{m-1} a_{ij}\, x^i y^j ,
\end{equation}
an element of the ring
$\mathbb{Z}_k[x,y]\,/\,(x^n-1,\ y^m-1)$: $x$-exponents are computed
modulo $n$, $y$-exponents modulo $m$, coefficients modulo $k$.  A
kernel translates into a polynomial $T(x,y)$ exactly as before; for
the kernel of Section~\ref{sec:kernel43},
\begin{equation}
\begin{aligned}
  T(x,y) ={}& -x^{-1}y^{-1} + 2x^{-1} - x^{-1}y
             + 2y^{-1} - 4 + 2y
             - xy^{-1} + 2x - xy\\
  \equiv{}& -x^{n-1}y^{m-1} + 2x^{n-1} - x^{n-1}y + 2y^{m-1} - 4 + 2y
            - xy^{m-1} + 2x - xy .
\end{aligned}
\end{equation}
Multiplying $A$ by $T$ yields the next generation.  Though the algebra
grows involved in higher dimensions, the polynomial representation is a
valuable conceptual tool: the superposition principle is simply the
distributive law of polynomial multiplication over addition.
Superposition says, in effect, that if you've seen one cell evolve,
you've seen all possible configurations of cells evolve.

Unfortunately, the most interesting CA are not additive and enjoy no
such superposition property.

\section{A Famous Non-Additive CA: The Game of Life}
\label{sec:gameoflife}

Conway's Game of Life is a simple 2-dimensional CA in which each cell
is either \emph{alive} or \emph{dead}.  Imagining the cells as squares
of a checkerboard, the neighborhood of a cell consists of its eight
neighboring squares (left, right, up, down, and the four diagonals).
The transition rules are:
\begin{itemize}
\item a dead cell becomes alive in the next generation exactly when it
  has exactly three living neighbors (a \emph{birth});
\item a live cell remains alive exactly when it has two or three
  living neighbors (otherwise it dies of loneliness or
  overcrowding).
\end{itemize}

In the Life figures below, white represents dead cells and live cells
are color-coded by age: cerise cells are newborn, red very young,
orange older, yellow older yet, green still older, blue very old, and
violet ancient --- more than 250 generations old.
Figure~\ref{fig:life2d} shows evolution from an initially random state;
after 500 generations the typical menagerie of Life structures has
emerged: stable \emph{still lifes} (blocks, beehives), oscillating
\emph{blinkers}, and the occasional \emph{glider} crawling across the
torus.  Figure~\ref{fig:liferow} shows evolution from a row of 71
initially live cells; the bilateral symmetry of the initial state is
preserved forever, and from generation 200 on, gliders can be seen at
each corner moving outward.

\begin{figure}[htbp]\centering
\includegraphics[width=0.48\textwidth]{ch4_life_initial}\;
\includegraphics[width=0.48\textwidth]{ch4_life_gen500}
\caption{Life from an initially random state with probability of life
$0.3$ (left); after 500 generations (right), various typical Life
structures have emerged, age-colored.}
\label{fig:life2d}
\end{figure}

\begin{figure}[htbp]\centering
\includegraphics[width=\textwidth]{ch4_life_row71}
\caption{Evolution from an initial row of 71 live cells.}
\label{fig:liferow}
\end{figure}

\expand{Life went on to a remarkable career.  In 1982 Conway and,
independently, Gosper's group established that Life is computationally
universal: glider streams can implement arbitrary logic circuits.
Self-constructing patterns were eventually built explicitly, notably
Gemini (2010), a pattern that constructs a complete copy of itself ---
achieving within Life the self-replication von Neumann originally
sought from cellular spaces.  Conway himself died in April 2020;
Life remains his most famous creation, to his occasional chagrin.}

\subsection*{A small bestiary}

The interest of Life lies in its \emph{structures}.  The simplest are
\emph{still lifes} such as the $2\times2$ block, frozen forever, and
simple oscillators such as the three-in-a-row \emph{blinker}, which
flips between horizontal and vertical with period 2.  Far more
consequential is the \emph{glider} (Figure~\ref{fig:glider}): a
five-cell pattern that, after four generations, reproduces its own
shape displaced one cell diagonally.  Gliders are Life's messengers ---
the moving particles out of which computation is built.

\begin{figure}[htbp]\centering
\includegraphics[width=\textwidth]{ch4_glider}
\caption{The glider returns to its original shape, translated by one
cell diagonally, every four generations --- a self-propelled pattern
moving across the torus.}
\label{fig:glider}
\end{figure}

That gliders can be manufactured to order was the key to Life's
universality.  Bill Gosper's \emph{glider gun}
(Figure~\ref{fig:gosper}) is a pattern of period 30 that emits a new
glider every 30 generations, a steady stream marching off toward
infinity.  Because a gun's population grows without bound, its
discovery in 1970 settled a \$50 wager posed by Conway --- whether any
Life pattern could grow forever --- and opened the door to building,
within Life, the logic gates, wires, and memory of a universal
computer.

\begin{figure}[htbp]\centering
\includegraphics[width=0.75\textwidth]{ch4_gosper_gun}
\caption{Gosper's glider gun after 120 generations: the period-30
mechanism at upper left has emitted a diagonal stream of four gliders.
The population grows without bound.}
\label{fig:gosper}
\end{figure}

The module \texttt{cadyn.ca2d} ships these patterns as data, ready to
place on a torus of any size:

\begin{lstlisting}
from cadyn import ca2d

# a glider in the corner of a 60x60 world, run for 40 generations
grid = ca2d.pattern_grid(ca2d.GLIDER, 60, top=2, left=2)
alive, age = ca2d.run_life(grid, 40)
ca2d.show_life(age)

# the Gosper gun; watch the population climb by 5 cells every 30 gens
gun = ca2d.pattern_grid(ca2d.GOSPER_GLIDER_GUN, 90, 130, top=5, left=5)
alive, age = ca2d.run_life(gun, 120)
\end{lstlisting}

\section{A Non-Additive CA Involving Euler's Totient Function}
\label{sec:totient}

Our second non-additive 2-dimensional example again takes $k=256$ on a
$500\times500$ torus, with the Moore neighborhood (the cell and its
eight nearest neighbors).  The transition rule borrows a function from
number theory, the Euler totient $\varphi$ (Appendix~\ref{app:totient}):
\begin{equation}
  C^{(t+1)}_{ij} \;=\; \sum_{a=-1}^{1}\sum_{b=-1}^{1}
     \varphi\bigl(C^{(t)}_{i+a,\,j+b}\bigr) \pmod{256},
\end{equation}
where we set $\varphi(0)=0$.  Starting from a single seed of value 1
(white $=$ 0), a mandala of intricate complexity and symmetry once more
unfolds; Figure~\ref{fig:totient} shows generation 250 with a
magnified detail.

\begin{figure}[htbp]\centering
\includegraphics[width=0.48\textwidth]{ch4_totient_gen250}\;
\includegraphics[width=0.48\textwidth]{ch4_totient_gen250_detail}
\caption{The Euler-totient rule at generation 250 (left) and a
magnified central detail (right).}
\label{fig:totient}
\end{figure}
